\scanpage{25}
% archive foliation: 24
\begin{center}- 95 -\end{center}
n° 268

variété de Picard, il nous suffira de montrer comment on peut définir de façon géométrique le passage au quotient de $A$ par le noyau de $\varphi_X$, tout en conservant certaines propriétés de rationalités, $X$ étant un diviseur non dégénéré.

Comme première application du théorème du carré, nous allons démontrer le plongement projectif d'une variété abélienne. Rappelons que si $X$ est un diviseur sur $A$ une variété $V$, complète et sans singularités de \unsure{c} codimension 1, on désigne par $\mathcal{L}(X)$ l'ensemble de tous les diviseurs positifs linéairement équivalents à $X$, et par $L(X)$ l'ensemble des fonctions $\varphi$ sur $V$ telles que $(\varphi)\geq-X$. C'est un espace vectoriel sur les constantes de dimension finie. Si $\mathcal{L}(X)$ n'est pas vide, on dit que c'est un système linéaire complet. (Cf. IAG). Si $V$ est définie sur $k$, et $X$ rationnel sur $k$, on peut trouver une base de $L(X)$ définie sur $k$.
On dit que $\mathcal{L}(X)$ \unsure{séparé|sépare} les points de $V$ si quels que soient les points $P,Q$ distincts de $V$, il existe un diviseur $Y\in\mathcal{L}(X)$ tel que $P\in\operatorname{supp}(Y)$ et $Q\notin\operatorname{supp}(Y)$. Si $\varphi_0,\ldots,\varphi_m$ est une base de $L(X)$, et $P$ est générique sur $V$ alors l'application $P\longrightarrow(\varphi_0(P),\ldots,\varphi_m(P))$ définit une application rationnelle de $V$ génériquement surjective sur une variété projective. Le système linéaire complet $\mathcal{L}$ est dit ample si l'application rationnelle est birationnelle et biholomorphe. L'existence d'un plongement projectif d'une variété normale $V$ équivaut à l'existence d'un système linéaire ample.

Ceci étant, nous avons le théorème suivant :

\uline{Théorème} 5 - Soit $A$ une variété abélienne. Alors il existe un plongement projectif de $A$, Si $X$ est un diviseur positif non dégénéré sur $A$, alors il existe un entier $m>0$ tel que $mX$ soit ample.
